\section{Experiments}
\label{sect-experiments}
We evaluate \texttt{Bach4Popper}, our coordination-based federated extension of Popper, on standard ILP benchmarks.
The objective of this experimental study is not to improve predictive performance over centralized ILP but to assess whether Popper's learning from failures mechanism can be executed in a coordinated, federated setting without compromising correctness and to quantify the symbolic cost induced by coordination and data decentralization. 

Centralized Popper serves as the primary baseline, representing the reference behavior when all examples and background knowledge are available on a single machine. Bach4Popper is evaluated under the same bias and search parameters, with the difference that examples are distributed across multiple clients and only abstract evaluation outcomes are exchanged through coordination primitives
\subsection{Experimental Goals}
Our evaluation is guided by the following research questions, in line with the theoritical results established in the methodology:
\begin{itemize}
    \item  \textbf{RQ1 (Correctness).}  Does \texttt{Bach4Popper} compute hypotheses that are equivalent to those learned by centralized Popper despite strict data locality constraints and the absence of shared examples?
    \item  \textbf{RQ2 (Performance and Coordination Cost).}  How does coordination-based federated execution impact the symbolic search process, in terms of refinement rounds and execution time, compared to centralized Popper?
%when increasing the number of participating clients from 2 to 5?
\end{itemize}

%To study scalability effects, we progressively increase the number of participating clients from two to five, while keeping the overall dataset and bias settings fixed.


%\subsection{Experimental Setup}
%We conduct experiments on three standard ILP benchmarks: 
%\texttt{Zendo}, \texttt{Trains}, and \texttt{IGGP-RPS}. 
%For each dataset, examples are partitioned across multiple clients to simulate a federated learning setting.
%Each client locally evaluates candidate hypotheses generated by the server and communicates only abstract evaluation outcomes through the Bach coordination language.

%Unless stated otherwise, experiments are conducted with two clients. We further extend the evaluation to three clients to assess the impact of increasing the number of participants on the coordination process.


\subsection{Datasets}
\label{sec:datasets}
We evaluate Bach4Popper on standard ILP benchmark datasets obtained from the official Popper implementation~\cite{popperv110}. These benchmarks are commonly used to analyze symbolic hypothesis refinement under controlled bias settings. Table~\ref{tab:datasets-config} summarizes the datasets and their corresponding Popper configurations.

\paragraph{Trains.}
The \texttt{Trains} dataset's goal is to induce rules differentiating eastbound from westbound trains~\cite{larson1977inductive}. Its a small dataset, but it demands relational reasoning over structured background knowledge and is frequently used to validate the correctness of ILP systems.

\paragraph{Zendo.}
\texttt{Zendo}'s dataset aims to infer a hidden concept from constructed structures~\cite{bramley2018grounding}, inspired by a rule-discovery game, used in relational concept learning benchmarks. 
It involves a larger hypothesis space and richer relational patterns  than trains; hence, it is suitable for evaluating symbolic generalization.

\paragraph{IGGP-RPS.}
The \texttt{IGGP-RPS} dataset originates from Inductive General Game Playing~\cite{cropper2020inductive} and targets the learning of rules that explain game-state transitions. This benchmark is challenging due to the size of the example set and the complexity of the relational background knowledge.

\begin{table*}[htb]
\centering
\caption{Datasets and corresponding Popper configuration used in the experiments.}
\label{tab:datasets-config}
\begin{tabular}{lcccccc}
\toprule
\textbf{Dataset} & \textbf{Target Pred.} & \textbf{\#Pos} & \textbf{\#Neg} & \textbf{MaxClauses} & \textbf{MaxVars} & \textbf{MaxBody} \\
\midrule
\texttt{Trains}   & \texttt{f/1}              & 5   & 5   & 3 & 5 & 6 \\
\texttt{Zendo1}   & \texttt{zendo/1}          & 20  & 20  & 4 & 5 & 7 \\
\texttt{iggp-rps} & \texttt{next\_score/1}    & 108 & 356 & 4 & 5 & 4 \\
\bottomrule
\end{tabular}
\end{table*}

\subsubsection{Federated Partitioning}
\label{sec:federated-partitioning}

We partition datasets into three client-specific partitions to simulate a cross-silo federated learning setting where the setting is composed of three clients and a central server. Each client holds a balanced subset of positive and negative examples and evaluates candidate hypotheses using only its local data. No examples, background knowledge, or intermediate evaluation results are shared between clients or with the server. The server have the bias data. 

Table~\ref{tab:fed-partitions} shows the different distributions of positive and negative examples across clients for each dataset. We use this partitioning across all federated experiments. 

\begin{table}[ht]
\centering
\caption{Positive and negative examples per dataset partition}
\label{tab:fed-partitions}
\setlength{\tabcolsep}{6pt}
\begin{tabular}{l|ccc|ccc|ccc}
\toprule
\textbf{Partition}
& \multicolumn{3}{c|}{\textbf{Trains}}
& \multicolumn{3}{c|}{\textbf{Zendo1}}
& \multicolumn{3}{c}{\textbf{IGGP-RPS}} \\
& POS & NEG & Total
& POS & NEG & Total
& POS & NEG & Total \\
\midrule

Part 1
& 2 & 2 & 4
& 7 & 7 & 14
& 36 & 119 & 155 \\

Part 2
& 2 & 2 & 4
& 7 & 7 & 14
& 36 & 119 & 155 \\

Part 3
& 1 & 1 & 2
& 6 & 6 & 12
& 36 & 118 & 154 \\

\bottomrule
\end{tabular}
\end{table}


\subsection{Results}


\subsubsection{Convergence.}
We first evaluate whether Bach4Popper preserves the descriptive correctness of centralized Popper.
To this end, we consider a hybrid evaluation protocol that combines federated learning with centralized testing.

In the federated setting, Bach4Popper is executed with three learning clients, each holding a disjoint partition of the dataset and evaluating hypotheses locally.
In addition, we introduce a centralized test server that has access to the full dataset and is used exclusively for evaluation purposes.
This test server does not participate in the learning or coordination process; it serves only to assess the final hypothesis induced by Bach4Popper on the complete dataset, ensuring comparability with centralized Popper.

For each dataset, we compare the hypothesis learned by centralized Popper with the hypothesis learned by \texttt{Bach4Popper} in terms of accuracy and recall.
To ensure a fair and consistent evaluation, we introduce a centralized \emph{test server} that has access to the full dataset and is used exclusively for evaluation purposes.
This test server does not participate in the learning or coordination process; instead, it evaluates the final hypothesis produced by \texttt{Bach4Popper} on the complete dataset, in the same manner as centralized Popper.
In parallel, the same hypothesis is also evaluated locally by each client on its own data partition.
This evaluation protocol allows us to assess whether a hypothesis learned through federated coordination exhibits equivalent behavior when tested on local data and on the full dataset.
As shown in Table~\ref{tab:convergence}, the coordinated learning process consistently converges to hypotheses that are complete and consistent on the full dataset, confirming that federation preserves the correctness guarantees of centralized Popper and answering \textit{RQ1}.



\begin{table*}[ht]
\centering
\small
\caption{Convergence of Bach4Popper towards centralized Popper in terms of descriptive performance.}
\label{tab:convergence}
\setlength{\tabcolsep}{6pt}
\begin{tabular}{l|cc|cc|cc}
\toprule
\textbf{Setting}
& \multicolumn{2}{c|}{\textbf{Trains}}
& \multicolumn{2}{c|}{\textbf{Zendo}}
& \multicolumn{2}{c}{\textbf{IGGP-RPS}} \\
& Accuracy & Recall
& Accuracy & Recall
& Accuracy & Recall \\
\midrule

Centralized
& 1.00  & 1.00 
& 1.00  & 1.00 
& 1.00  & 1.00   \\

Bach4Popper
& 1.00  & 1.00 
& 1.00  & 1.00 
& 1.00  & 1.00   \\
\bottomrule
\end{tabular}
\end{table*}

%We report the hypotheses learned across three independent runs for each dataset in Table~\ref{tab:hypotheses}.
%For each run, the hypothesis induced by Bach4Popper is logically equivalent to the one learned by centralized Popper, although their syntactic structure may differ due to the non-deterministic nature of the symbolic search.
These experimental results provide an empirical validation of Propositions~1-4 established in Section~\ref{sect-BachForPopper}, showing that symbolic outcomes and scores computed in a federated manner induce the same hypothesis refinement behavior as centralized Popper.

\subsubsection{Hypothesis Variability}
Although correctness is preserved, the concrete hypotheses learned across runs may differ syntactically.
This variability is already present in centralized Popper due to its heuristic-driven ASP search and is further amplified in the federated setting by data partitioning.
%Table~\ref{tab:hypotheses} illustrates that different executions may converge to distinct but logically equivalent hypotheses. These differences reflect alternative valid solutions within the same hypothesis space rather than learning errors.
Table~\ref{tab:hypotheses-iggp} shows that, although Bach4Popper may converge to syntactically different hypotheses across runs, the induced programs remain
logically equivalent to the hypothesis learned by centralized Popper. Each hypothesis consists of multiple rules and satisfies completeness and
consistency under the same language bias.

% was in caption: Bach4Popper may converge to different multi-rule hypotheses across runs,all of which are logically equivalent to the centralized Popper solution.

\begin{table}[ht]
\centering
\scriptsize
\caption{Illustration of hypothesis variability for \texttt{IGGP-RPS} across two independent federated runs and the centralized Popper solution.}

\label{tab:hypotheses-iggp}
\setlength{\tabcolsep}{6pt}
\begin{tabular}{l|p{0.78\linewidth}}
\toprule
\textbf{Setting} & \textbf{Learned hypothesis} \\
\midrule

Federated (Run 1) &
\texttt{next\_score(A,B,C) :- different(D,E),different(E,B), my\_true\_score(A,D,C),different(D,B).} \\
& \texttt{next\_score(A,B,C) :- different(D,E),my\_true\_score(A,D,C), different(B,E), wins(A,E).} \\
& \texttt{next\_score(A,B,C) :- player(B),my\_succ(D,C), wins(A,B), my\_true\_score(A,B,D).} \\
& \texttt{next\_score(A,B,C) :- does(A,B,E),does(A,D,E), different(B,D),
my\_true\_score(A,B,C).}
\\
\midrule

Federated (Run 2) &
\texttt{next\_score(A,B,C) :- my\_succ(E,C), my\_true\_score(A,B,E), my\_true\_score(A,D,C), my\_true\_score(A,D,E).} \\
& \texttt{next\_score(A,B,C) :- wins(A,E), different(E,B), my\_true\_score(A,D,C), different(D,E).} \\
& \texttt{next\_score(A,B,C) :- my\_true\_score(A,B,D), player(B), wins(A,B), my\_succ(D,C).} \\
& \texttt{next\_score(A,B,C) :- different(D,B), does(A,D,E), my\_true\_score(A,B,C), does(A,B,E).}
\\
\midrule

Centralized Popper &
\texttt{next\_score(A,B,C) :- my\_succ(D,C), wins(A,B), my\_true\_score(A,B,D).} \\
& \texttt{next\_score(A,B,C) :- different(D,B), wins(A,D), my\_true\_score(A,B,C).} \\
& \texttt{next\_score(A,B,C) :- different(B,E), does(A,B,D), my\_true\_score(A,B,C), does(A,E,D).}
\\

\bottomrule
\end{tabular}
\end{table}




\subsubsection{Performance and Symbolic Cost}
We compare the symbolic cost of centralized and federated execution by measuring the number of refinement rounds and the total execution time required to reach termination.
To account for the non-deterministic nature of Popper’s search, all performance results in Table~\ref{tab:perf} are averaged over ten independent runs for each dataset and setting.
For each run, we record the number of refinement rounds and the execution time measured from the generation of the first hypothesis until termination.
Reported values correspond to the arithmetic mean across successful runs. Indeed, we recall that an execution terminates under one of three conditions:
\begin{enumerate}
    \item a globally valid hypothesis $(\texttt{all}, \texttt{none})$ is found,
\item the hypothesis space is exhausted under the imposed bias constraints, or
\item a predefined timeout is reached.
These termination conditions mirror those of centralized Popper and are preserved in the federated setting.
\end{enumerate}
%Due to the heuristic-driven nature of the symbolic search, not all runs necessarily converge to a globally valid hypothesis within the allotted resources.%In particular, convergence depends on the trajectory of the search and on the initial hypotheses explored. When a run does not reach a solution before the search space is exhausted or the timeout is reached, the experiment is restarted and only successful runs are retained for averaging.

\begin{table*}[ht]
\centering
\small
\caption{Performance comparison for centralized Popper vs Bach4Popper, averaged over X independent runs.}
\label{tab:perf}
\setlength{\tabcolsep}{6pt}
\begin{tabular}{l|cc|cc|cc}
\toprule
\textbf{Setting}
& \multicolumn{2}{c|}{\textbf{Trains}}
& \multicolumn{2}{c|}{\textbf{Zendo}}
& \multicolumn{2}{c}{\textbf{IGGP-RPS}} \\
& \#Rounds & Time (s)
& \#Rounds & Time (s)
& \#Rounds & Time (s) \\
\midrule

Centralized
& 105 & 0,67
& 185 & 0,79
& 1835 & 11,99\\

\midrule

Bach4Popper
& 2178 &58,75
& 5998 &202,40
& 11080 & 350 \\ 

\bottomrule
\end{tabular}
\end{table*}

Federated execution generally incurs a higher number of refinement rounds and longer execution times than centralized Popper, reflecting the additional coordination and synchronization overhead.
However, this overhead is not systematic.
In some runs, particularly for \textit{IGGP-RPS}, federated execution converges in a comparable or even smaller number of rounds than centralized execution, answering \textit{RQ2}.

\subsection{Analysis and Discussion}
\label{sec:discussion}

\paragraph{Correctness.}
Across all benchmarks, Bach4Popper preserves the correctness guarantees of centralized Popper.
Despite strict data locality constraints and the absence of shared examples or background knowledge, the coordinated learning process consistently converges to hypotheses that are complete and consistent under the same language bias.
As shown in Table~\ref{tab:convergence}, centralized and federated executions achieve identical descriptive performance, confirming that coordination does not affect logical soundness which answers to \textit{RQ1}.

\paragraph{Hypothesis variability.}
Although correctness is preserved, the concrete hypotheses learned across runs may differ syntactically.
This variability is already present in centralized Popper due to its heuristic-driven ASP search, and it is further amplified in the federated setting by data partitioning.
As illustrated in Table~\ref{tab:hypotheses}, different executions may converge to distinct but logically equivalent hypotheses.
These differences reflect alternative valid solutions within the same hypothesis space rather than learning errors, and are typically associated with different ASP termination paths during the search.

\paragraph{Cost of federation.}
Federation impacts the efficiency of the symbolic search but not its outcome. Because clients evaluate hypotheses on partial data, some hypotheses may appear locally plausible while being globally inconsistent, requiring additional refinement rounds.
This behavior explains the increased symbolic cost observed in federated execution.

%Federation impacts the efficiency of the symbolic search but not its outcome. Because each client evaluates candidate hypotheses on partial data, some hypotheses may appear locally plausible while being globally inconsistent. As a consequence, additional refinement rounds may be required to reach a globally acceptable hypothesis. This behavior is reflected in Table~\ref{tab:perf}, where federated execution often incurs more rounds and longer execution times than centralized Popper. Importantly, this overhead is not systematic. In some runs, for \textit{iggp-rps},  coordination leads to convergence in a comparable or even smaller number of rounds than centralized execution, answering \textit{RQ2}.

%The observed variability across runs is partly explained by the dependence on the initial hypotheses explored during the search.
%Since Popper relies on heuristic-driven ASP solving, different initial models may lead to distinct refinement trajectories.
%This effect is particularly pronounced in complex relational domains such as \textit{IGGP-RPS}, where early hypotheses often partially cover both positive and negative examples, resulting in longer convergence times or early exhaustion in some runs.

\paragraph{Limitations.}
Our evaluation focuses on controlled ILP benchmarks and does not aim to demonstrate large-scale federated deployment. Exploring larger and noisier relational domains, heterogeneous client distributions remains an important direction for future work.

